% ============================================================
\section{Ejercicio 2}

% ============================================================

\begin{tcolorbox}[enunciado]
Considera la demostración vista en clase del Teorema de Cook-Levin.
\begin{enumerate}
    \item ¿La demostración podría funcionar si definimos las ventanas como rectángulos de 3 renglones y dos columnas? ¿Por qué?
    \item ¿Sería posible usar ventanas más grandes para realizar la demostración? por ejemplo de 3 x 3 ¿Por qué?
    \item ¿Por qué son necesarias cada una de las subfórmulas en la transformación propuesta en la demostración del teorema? ¿Qué pasaría si quitamos una de ellas?
    \item ¿Cómo cambiaría la demostración de la NP-Dificultad si utilizamos la definición alternativa de NP? es decir, en vez de suponer la existencia de una máquina No Determinista, suponemos la existencia de una máquina verificadora.
    \item Considera alguna de las máquinas de Turing de las tareas anteriores.
    \begin{itemize}
        \item Ilustra ejemplos de ventanas legales y ventanas ilegales. Ilustra y especifica las fórmulas booleanas correspondientes a las ventanas válidas.
        \item Ilustra ejemplos de algunas de las subfórmulas para un cómputo concreto en donde se acepta la cadena de entrada.
        \item Ilustra ejemplos de algunas de las subfórmulas para un cómputo concreto en donde NO acepta la cadena de entrada.
    \end{itemize}
    Justifica ampliamente tus respuestas.
\end{enumerate}
\end{tcolorbox}

% ------------------------- Solución -------------------------
\section{Respuestas}
\subsection{1}
    No, ya que el teorema establece ventanas validas de 2 renglones x 3 columnas, pues la codificación de las configuraciones podría cambiar en 3 posiciones consecutivas con una sola transición. principalmente tenemos que\\ 
    
    $P_{i,j} = P_{celda} \land P_{incio} \land P_{aceptacion} \land P_{mov}$\\
    
    Para cada fila se tiene el estado siguiente, el simbolo actual, simbolo de remplazo en la cinta. Si quitamos la primera o ultima columna para hacer una ventana valida de 3 x 2 estariamos perdiendo informacion con respecto a las posibles configuraciones posbiles que se pueden realizar en $P_{mov}$.
    
    Es decir que a partir de 2 movimientos se pueda determinar un tabluex satisfacible en esa venta de 3 x 2 sin determinar si se realizo un cambio de simbolo al final de esta transicion o se mantuvo, lo que provocaria una inconstistencia al momento de analizar cada ventana valida.

    Pues en un solo paso, el cabezal de una máquina de Turing solo puede desplazarse una celda hacia la izquierda o una celda hacia la derecha. Las 3 columnas constituye el mínimo necesario para registrar el movimiento del cabezal.
\subsection{2}
Sí, pues al agregar una tercera fila es decir, evaluar dos transiciones consecutivas de la Máquina de Turing sobre 3 columnas contiguas se conserva toda la información  necesaria sin perder los bordes laterales, lo que mantiene el tamaño polinomial al verificar la satisfacibilidad dentro de las ventanas válidas en $P_{\mathrm{mov}} = \bigwedge_{1 \le i,j \le n^k - 2} (\text{ventana válida en } (i,j))$. Por lo tanto, solo se tendría que obtener todas las ventanas válidas posibles de $3 \times 3$ que cumplen con las reglas de transición hasta la tercera fila la cual seria la configuración $i+2$. De esta manera, la fórmula se establecería de la siguiente manera:

    $$\begin{aligned} P_{\mathrm{mov}} = \bigwedge_{1 \le i,j \le n^k - 2} \Big\{ \, & (x_{i,j,a_1} \,\wedge\, x_{i,j+1,a_2} \,\wedge\, x_{i,j+2,a_3} \\ & \wedge\, x_{i+1,j,a_4} \,\wedge\, x_{i+1,j+1,a_5} \,\wedge\, x_{i+1,j+2,a_6} \\ & \wedge\, x_{i+2,j,a_7} \,\wedge\, x_{i+2,j+1,a_8} \,\wedge\, x_{i+2,j+2,a_9}) \\ & \vee \\ & (x_{i,j,b_1} \,\wedge\, x_{i,j+1,b_2} \,\wedge\, x_{i,j+2,b_3} \\ & \wedge\, x_{i+1,j,b_4} \,\wedge\, x_{i+1,j+1,b_5} \,\wedge\, x_{i+1,j+2,b_6} \\ & \wedge\, x_{i+2,j,b_7} \,\wedge\, x_{i+2,j+1,b_8} \,\wedge\, x_{i+2,j+2,b_9}) \\ & \vee\, \cdots \,\Big\} \end{aligned}$$

\subsection{3}

   Porque cada subfórmula garantiza una condición estructural indispensable del cuadro de cómputo (tableau), de forma que la fórmula global $P = P_{\text{celda}} \land P_{\text{inicio}} \land P_{\text{aceptación}} \land P_{\text{mov}}$ sea satisfacible si y solo si existe un cómputo válido de la Máquina de Turing $N$ que acepta la cadena $w$.\\

    Tenemos que $P_{\text{mov}} = \bigwedge_{1 \le i,j \le n^k - 2} (\text{ventana válida en } (i,j))$, la cual verifica que el paso de cada fila (la cual son las configuraciones) a la siguiente respete las transiciones legales definidas por $\delta$ en $N$.\\
    
    Para $P_{\text{celda}} = \bigwedge_{1 \le i,j \le n^k} \left[ \left(\bigvee_{s \in C} x_{i,j,s}\right) \land \left(\bigwedge_{s \neq t} (\neg x_{i,j,s} \lor \neg x_{i,j,t})\right) \right]$, esta subfórmula nos garantiza que toda celda tiene al menos un símbolo y a lo sumo un símbolo del alfabeto del tableau $C$, es decir, exactamente un símbolo por celda.\\

    Al igual que $P_{inicio}$, muestra el símbolo exacto $s$ que se encuentra en cada posición $(1,j)$ de la primera fila;, esto quiere decir que la primera fila de la tabla es igual a la configuración inicial de N en w..\\

        $$
        \begin{aligned}
        P_{\mathrm{inicio}} =\;&
        x_{1,1,\#}
        \;\wedge\;
        x_{1,2,q_0}
        \;\wedge\;
        x_{1,3,w_1}
        \;\wedge\;
        x_{1,4,w_2}
        \;\wedge\; \cdots \\
        &\wedge\;
        x_{1,n+2,w_n}
        \;\wedge\;
        x_{1,n+3,\square}
        \;\wedge\; \cdots \\
        &\wedge\;
        x_{1,n^k-1,\square}
        \;\wedge\;
        x_{1,n^k,\#}.
        \end{aligned}
        $$

    Finalmente, la última subfórmula la cual es $P_{\text{aceptación}}$ nos permite verificar que se alcanza el estado de aceptación si al menos una celda del cuadro contiene el estado $q_a$, expresada como $\bigvee_{1 \leq i,j \leq n^k} x_{i,j,q_a}$; de esta manera podemos saber si la cadena fue aceptada o no por dicha MT.\\

    Cada subformula corresponde a una restricción necesaria que debe cumplirse simultáneamente, ya que, como se debe de verificar que exista una tautologia, no se podría colocar unificadamente en una sola formula cada regla o existirian inconsistencias, por ejemplo, si para la formula $P_{celda}$ se tomara la disyuncion $\bigvee$, para unificar $P_{mov}$, entonces podrian existir negaciones de las clausulas, con sus contrapositivas, lo que ocasionaría que se cancelasen, lo cual, alteraria el resultado de la formula real y el computo seria invalido.\\
    
    Ahora, si quitáramos alguna de las reglas, como por ejemplo $P_{\text{celda}}$, una celda podría quedar vacía o contener múltiples símbolos a la vez, ahora si quitáramos $P_{\text{inicio}}$, el cómputo podría comenzar en una configuración arbitraria distinta de $w$, por lo que con esto podría empezar directamente en $q_a$; si quitáramos $P_{\text{mov}}$, las filas podrían saltar  entre estados sin seguir la función de transición $\delta$, finalmente si quitáramos $P_{\text{aceptación}}$, cualquier cómputo válido que rechace o no termine también satisfaría la fórmula, haciendo que el resultado fuera inválido.\\

    Por lo tanto las subfórmulas codifican las 4 condiciones necesarias para que exista un cuadro de aceptación para el cómputo de $N$ con la entrada $w$.

\subsection{4}
    Por definición de la clase NP, existen una función polinomial $p$ y una máquina de Turing $M$ que opera en tiempo polinomial, tales que:
$x \in L \iff \exists u \in \{0, 1\}^{p(|x|)}$ tal que $M(x, u) = 1$.\\

Para demostrar que SAT es NP-dificil, debemos presentar una función $f$, computable en tiempo polinómico, que transforme cualquier $x \in \{0, 1\}^*$ en una fórmula en forma normal conjuntiva (CNF) $\omega_x$, de modo que $x \in L$ si y solo si $\omega_x$ es satisfacible. Entonces de esta manera se puede fijar una entrada $x$ de longitud $n$, un certificado $u$ de longitud $p(n)$, y definimos $y = x \circ u$, cuya longitud es $n + p(n)$.\\

Se propondrá un conjunto de variables booleanas (que codifican los bits de $y$ y el estado de la máquina en cada paso), a las cuales se les atribuirán 2 posibles valores, verdadero ($1$) o falso ($0$), estructuradas en un conjunto de cláusulas.\\

Se asume que $M$ es una máquina de dos cintas una de entrada de solo lectura para $y$ y otra de trabajo, y hay que tomar tambien en cuenta la posición de sus cabezales en el paso $i$ que depende únicamente de $i$ y de la longitud de la entrada, no del contenido de las cintas. Cada paso $i \in \{0, \dots, T(n)\}$ se representa mediante una 'captura' $z_i \in \{0,1\}^c$ que codifica el estado actual y los símbolos leídos bajo ambos cabezales, donde $c = \vert{}Q\vert{} \times \vert{}\Gamma\vert{} \times \vert{}\Gamma\vert{}$ es una constante independiente de $n$. \\

\textbf{Construcción}

La fórmula $\omega_x(y, z_0, \dots, z_{T(n)})$ se construye como la conjunción de cuatro subfórmulas en CNF:\\ 

Fijación de $x$ y libertad de $u$: Cláusulas que exigen que los primeros $n$ bits de $y$ coincidan exactamente con la entrada dada $x$. Los $p(n)$ bits restantes ($y_{n+1}, \dots, y_{n+p(n)}$) se dejan completamente libres para que cualquier asignación de verdad actúe como un certificado candidato $u$.\\

Configuración inicial ($z_0$): Cláusulas que obligan a que $z_0$ codifique el estado inicial $q_0$ con los cabezales en su posición de arranque. La cadena $z_0$ codifica la 'captura' inicial $< s, \vdash,  \vdash>$ de M.\\

Transiciones locales deterministas ($z_i$): Dado que $M$ es determinista y modifica a lo sumo un símbolo por paso en la cinta de trabajo, cada instantánea $z_i$ depende únicamente de la instantánea inmediata anterior $z_{i-1}$, de la última instantánea $z_{\text{prev}(i)}$ en que el cabezal de trabajo estuvo en esa misma celda, y del bit leído en la cinta de entrada $y_{\text{pos}(i)}$. Esto se expresa mediante una función determinista $z_i = F(z_{i-1}, z_{\text{prev}(i)}, y_{\text{pos}(i)})$ para todo $1 \le i \le T(n)$.\\

Condición de aceptación ($z_{T(n)}$): Cláusulas que exigen que la última 'captura' $z_{T(n)}$ corresponda a un estado de parada donde $M$ produce como salida $1$.\\

\textbf{Complejidad en Tiempo}

Como mencionamos anteriormente convertir directamente la función booleana $u \mapsto M(x,u)$ a CNF produciría una fórmula de tamaño exponencial $p(n)2^{p(n)}$.\\ 

Pero si descomponemos el cómputo en $O(n)$ en pasos locales lo reducimos a manera polinomial. \\   

Al fijar los primeros $n$ bits de $y$ requiere $4n$ símbolos.  Las condiciones de inicio $z_0$ y aceptación $z_{T(n)}$ requieren tamaño constante $c2^c$ cada una.\\

Dado que cada paso de transición $z_i = F(\dots)$ involucra solo $3c + 1$ variables booleanas, esto hace que su tabla de verdad convertida a CNF ocupa un tamaño constante $(3c+1)2^{3c+1}$. \\

En conjunto, $\omega_x$ tiene tamaño $d \cdot (n + O(n))$ para una constante $d$ que depende solo de $M$. Dado que $O(n)$ es un polinomio en $n$, la fórmula tiene tamaño polinomial y el algoritmo de reducción $f$ la genera en tiempo polinomial sin necesidad de evaluar las $2^{p(n)}$ posibles asignaciones de $u$

    Comprobamos la equivalencia. ($\Rightarrow$): Si $x \in L$, existe un certificado $u \in \{0,1\}^{p(n)}$ tal que $M(x,u) = 1$. Al asignar $y = x \circ u$ y dar a $z_0, \dots, z_{T(n)}$ los valores exactos de las instantáneas generadas por $M$ durante esa ejecución, se satisfacen las cuatro condiciones y $\omega_x$ resulta satisfacible.
    
    Ahora para la otra parte  ($\Leftarrow$) si $\omega_x$ es satisfacible, existe una asignación $(y, z_0, \dots, z_{T(n)})$ que hace verdaderas todas sus cláusulas. Por la primera condición, $y$ tiene la forma $x \circ u$ para algún $u \in \{0,1\}^{p(n)}$; por la segunda y tercera, la secuencia $z_i$ corresponde fielmente a la única ejecución por parte de  $M$ sobre $(x, u)$, finalmente la ultima condición, tal   ejecución termina aceptando $M(x,u) = 1$, lo que implica $x \in L$.

    Por lo tanto $SAT$ es $NP-Dificil$.
    \cite{SinhaChaturvedi2014}
\subsection{5}
\subsubsection{1}

\subsubsection{2}

\subsubsection{3}
